\documentclass[12pt]{article}
\usepackage[utf8]{inputenc}
\usepackage{amsmath, amssymb, amsthm}
\usepackage[margin=1in]{geometry}
\usepackage{algorithm}
\usepackage{algorithmic}

\newtheorem{theorem}{Theorem}
\newtheorem{lemma}[theorem]{Lemma}

\title{Problem 85: Counting Rectangles}
\author{Project Euler Solution}
\date{}

\begin{document}
\maketitle

\section{Problem Statement}

Find the area of the rectangular grid whose number of contained axis-aligned rectangles is nearest to $T = 2{,}000{,}000$.

\section{Mathematical Foundation}

\begin{theorem}[Rectangle Counting Formula]
An $m \times n$ grid contains exactly
\[
R(m, n) = \binom{m+1}{2} \binom{n+1}{2} = \frac{m(m+1)}{2} \cdot \frac{n(n+1)}{2}
\]
axis-aligned rectangles.
\end{theorem}

\begin{proof}
An $m \times n$ grid is delineated by $m+1$ vertical lines and $n+1$ horizontal lines. An axis-aligned rectangle is uniquely determined by choosing 2 of the $m+1$ vertical lines (left and right edges) and 2 of the $n+1$ horizontal lines (top and bottom edges). These choices are independent, giving $\binom{m+1}{2} \cdot \binom{n+1}{2}$.
\end{proof}

\textbf{Verification:} For $m = 3, n = 2$: $R(3,2) = 6 \cdot 3 = 18$. $\checkmark$

\begin{lemma}[Search Bound]
We need only search $m$ values satisfying $\frac{m(m+1)}{2} \leq T$, i.e., $m = O(\sqrt{T})$.
\end{lemma}

\begin{proof}
Since $R(m, n) \geq R(m, 1) = \frac{m(m+1)}{2}$ for all $n \geq 1$, once $\frac{m(m+1)}{2} > T$ the rectangle count exceeds the target for every valid $n$. The quadratic inequality $m(m+1)/2 \le T$ yields $m \le \lfloor(-1 + \sqrt{1 + 8T})/2\rfloor = O(\sqrt{T})$.
\end{proof}

\begin{lemma}[Optimal $n$ for Fixed $m$]
For fixed $m$, the real root $n^*$ of $R(m, n^*) = T$ is
\[
n^* = \frac{-1 + \sqrt{1 + 8T/[m(m+1)]}}{2}.
\]
The optimal integer $n$ is either $\lfloor n^* \rfloor$ or $\lceil n^* \rceil$, since $R(m,n)$ is strictly increasing in~$n$.
\end{lemma}

\begin{proof}
Setting $R(m,n) = T$ gives $n(n+1) = 4T/[m(m+1)]$. The quadratic formula yields the stated positive root. Strict monotonicity of $R$ in $n$ ensures the closest integer value lies at the floor or ceiling.
\end{proof}

\section{Algorithm}

\begin{algorithm}[H]
\caption{Find Grid with Nearest Rectangle Count to $T$}
\begin{algorithmic}[1]
\STATE $\text{best\_diff} \gets \infty$, $\text{best\_area} \gets 0$
\FOR{$m = 1, 2, 3, \ldots$}
    \STATE $T_m \gets m(m+1)/2$
    \IF{$T_m > T$}
        \STATE \textbf{break}
    \ENDIF
    \STATE $n^* \gets \frac{-1 + \sqrt{1 + 8T/[m(m+1)]}}{2}$
    \FOR{$n \in \{\lfloor n^* \rfloor, \lceil n^* \rceil\}$}
        \IF{$n \geq 1$}
            \STATE $R \gets T_m \cdot n(n+1)/2$
            \IF{$|R - T| < \text{best\_diff}$}
                \STATE Update $\text{best\_diff}$ and $\text{best\_area} \gets m \cdot n$
            \ENDIF
        \ENDIF
    \ENDFOR
\ENDFOR
\RETURN $\text{best\_area}$
\end{algorithmic}
\end{algorithm}

\section{Complexity Analysis}

\begin{itemize}
    \item \textbf{Time:} $O(\sqrt{T})$. The loop runs while $m(m+1)/2 \leq T$, so $m = O(\sqrt{T})$. Each iteration does $O(1)$ work.
    \item \textbf{Space:} $O(1)$.
\end{itemize}

\section{Answer}

\[
\boxed{2772}
\]

\end{document}
